\documentclass[12pt]{article}
\usepackage[margin=1in]{geometry}
\usepackage{amsmath, amssymb, amsfonts}
\usepackage{graphicx}
\usepackage{siunitx}

\title{ELCT 562 PA6}
\author{Jacob Vaught}
\date{\today}

\begin{document}

\maketitle

\section*{Problem \#1}
Consider a set of signals defined by
\[
s_i(t) = \cos(2\pi f_c t + \theta_n), \quad nT_s \leq t \leq (n + 1)T_s
\]
where $i \in \{0,1,\ldots, M-1\}$, and $\theta_n \in \frac{2\pi i}{M}$. A communication system used by a platform on Mars employs these signals with $M = 2$ and operates at a data rate $R_b = 2\, \text{kbps}$ over a satellite link. The system parameters are as follows:
\begin{itemize}
    \item Distance between transmitter and receiver: $56 \times 10^6\, \text{km}$
    \item Carrier frequency: $1\, \text{GHz}$
    \item Transmitter power: $100\, \text{milliwatts}$
    \item Transmitter antenna gain: $47\, \text{dBi}$
    \item Receiver antenna gain: $30\, \text{dBi}$
    \item Receiver noise figure: $2\, \text{dB}$
\end{itemize}
The task is to analyze this communication system by answering the following questions:
\begin{enumerate}
    \item What kind of modulation (or signaling) does this communication system use?
    \item What is the pulse shape utilized for this communication system?
    \item What is the value of $T_s$?
    \item What is the received signal power in dBm based on the free-space path loss model?
    \item What is the energy per data symbol ($E_s$) in dBm at the receiver?
    \item What is the energy per bit ($E_b$) in dBm at the receiver?
    \item What is the complex noise variance ($N_0$) in dBm at the receiver?
    \item Can we achieve a bit-error rate of $10^{-2}$? Why?
    \item If not, what is the bit rate we can use to obtain a $0.01$ error probability when $E_b/N_0 = 4.1\, \text{dB}$?
\end{enumerate}

\section*{Solution}
\subsection*{Part (a): Type of Modulation}
The system uses Binary Phase Shift Keying (BPSK) as the modulation scheme, evident from the phase shifts representing two distinct symbols.

\subsection*{Part (b): Pulse Shape}
Typically, for BPSK systems, the pulse shape in baseband would correspond to a rectangular pulse or non-return-to-zero (NRZ) signaling.

\subsection*{Part (c): Value of $T_s$}
Given the data rate $R_b = 2\, \text{kbps}$, the symbol duration $T_s$ is calculated as:
\[
T_s = \frac{1}{R_b} = \frac{1}{2000} = 0.0005\, \text{seconds} = 500\, \text{microseconds}.
\]

\subsection*{Part (d): Received Signal Power}
The received signal power \( P_r \) can be calculated using the Friis transmission equation in decibels:
\[
P_r = P_t + G_t + G_r - L_{fs}
\]
where:
\begin{itemize}
    \item \( P_t = 10 \log_{10}(100 \times 10^{-3}) \) dBm is the transmitter power in dBm.
    \item \( G_t = 47 \) dBi is the transmitter antenna gain.
    \item \( G_r = 30 \) dBi is the receiver antenna gain.
    \item \( L_{fs} = 20 \log_{10}(d) + 20 \log_{10}(f) + 20 \log_{10}(\frac{4\pi}{c}) \) dB is the free-space path loss, with \( d = 56 \times 10^6 \) km converted to meters, \( f = 1 \) GHz, and \( c = 3 \times 10^8 \) m/s (speed of light).
\end{itemize}
Plugging in the values and calculating \( P_r \), we find that the received signal power \( P_r \) is approximately \(-150.41\) dBm.

\subsection*{Part (e): Energy per Data Symbol}
The energy per data symbol \( E_s \) at the receiver can be calculated using the formula:
\[
E_s = P_r + 10 \log_{10}(T_s)
\]
where \( T_s = \frac{1}{R_b} \) is the symbol duration. For \( R_b = 2 \) kbps, \( T_s = 0.0005 \) seconds (or 500 microseconds). Plugging in the values for \( P_r \) and \( T_s \), we find that the energy per data symbol \( E_s \) is approximately \(-183.42\) dBm.

\subsection*{Part (f): Energy per Bit}
For BPSK modulation, each symbol carries one bit of information. Therefore, the energy per bit \( E_b \) is equal to the energy per symbol \( E_s \):
\[
E_b = E_s
\]
Thus, the energy per bit \( E_b \) is also approximately \(-183.42\) dBm.

\subsection*{Part (g): Complex Noise Variance}
The complex noise variance \( N_0 \) at the receiver can be calculated using:
\[
N_0 = kT + NF
\]
where \( kT = -174 \) dBm/Hz represents the thermal noise power density at room temperature, and \( NF = 2 \) dB is the noise figure of the receiver. Adding these values, the complex noise variance \( N_0 \) at the receiver is calculated to be \(-172\) dBm/Hz.

\textbf{Note: This may be incorrect. It does not take into account frequency, because it makes the solution almost impossible to actually use}



\subsection*{Part (h): Bit-Error Rate}
The BER for BPSK in AWGN is given by:
\[
P_e = Q\left(\sqrt{2 \frac{E_b}{N_0}}\right)
\]
For the calculated \(E_b/N_0\) ratio, this results in:
\[
P_e = Q\left(\sqrt{2 \times \text{[linear value of } E_b/N_0]}\right) \approx 0.352
\]
indicating a BER higher than the target \(10^{-2}\), hence not feasible under current parameters.

\subsection*{Part (i): Required Bit Rate for \(E_b/N_0 = 4.1\) dB}
To achieve a BER of \(10^{-2}\), the required bit rate \(R_b'\) is adjusted based on:
\[
R_b' = R_b \times \frac{\text{[current } E_b/N_0 \text{ in linear]}}{\text{[required } E_b/N_0 \text{ in linear]}} \approx 56.16 \text{ bps}
\]
This new bit rate ensures the system meets the BER requirement of \(10^{-2}\) for \(E_b/N_0 = 4.1\) dB.


\section*{Problem\#2}

Consider a passband signal with the following characteristics:
\begin{itemize}
    \item Center frequency: $2.45 \text{ GHz}$
    \item Frequency range: $|f| \leq 10 \text{ MHz}$ around the center frequency.
\end{itemize}
\noindent
Determine:
\begin{enumerate}
    \item[(a)] The total bandwidth required for the signal.
    \item[(b)] The minimum sample rates on I and Q branches for a complex baseband representation.
    \item[(c)] The maximum data rate achievable using BPSK modulation.
\end{enumerate}

\subsection*{Solution}

\subsubsection*{(a) Bandwidth Allocation:}

The bandwidth required for the passband signal, which is centered around $2.45 \text{ GHz}$ with a frequency range of $|f| \leq 10 \text{ MHz}$, spans from $2.44 \text{ GHz}$ to $2.46 \text{ GHz}$. This means the total bandwidth needed for this signal is:

\[
\text{Total Bandwidth} = (2.46 \text{ GHz} - 2.44 \text{ GHz}) = 20 \text{ MHz}.
\]

\subsubsection*{(b) Minimum Sample Rates on I and Q Branches:}

For a complex (I/Q) baseband representation, each branch must accurately represent the full bandwidth of the signal to ensure all the necessary information is captured. Since the total bandwidth of the baseband signal is $20 \text{ MHz}$, covering from $-10 \text{ MHz}$ to $+10 \text{ MHz}$ after downconversion, the Nyquist theorem dictates that the minimum sampling rate required for each of the I and Q branches must be at least twice the highest frequency present in the baseband signal to prevent aliasing. Hence, each branch should be sampled at:

\[
\text{Minimum Sample Rate} = 2 \times 10 \text{ MHz} = 20 \text{ MHz}.
\]

\subsubsection*{(c) Maximum Data Rate for BPSK:}

For Binary Phase Shift Keying (BPSK), the maximum data rate is typically equal to the bandwidth of the signal, as each symbol represents one bit. Assuming ideal conditions (ignoring factors such as excess bandwidth for filtering), the maximum bit rate for BPSK, given the bandwidth of the baseband signal is $10 \text{ MHz}$, would be:

\[
\text{Maximum Data Rate} = 10 \text{ Mbps}.
\]

This analysis assumes ideal conditions without considering real-world factors such as channel coding, inter-symbol interference, or other forms of signal degradation.
\subsubsection*{(d) Sketches:}

The transformation of a signal from baseband to passband is a crucial concept in communications systems. Below, we illustrate this process through the frequency spectra of baseband and passband signals.

\begin{figure}[h!]
    \centering
    \includegraphics[width=0.7\textwidth]{a.png}
    \caption{Baseband (I \& Q branches)}
    \label{fig:baseband}
\end{figure}

Figure \ref{fig:baseband} shows the frequency spectra of the baseband signals for both the In-phase (I) and Quadrature (Q) components. The baseband signals are represented within the frequency range of $\pm10$ MHz. The solid line corresponds to the I branch, and the dashed line corresponds to the Q branch. These spectra are ideal representations, displaying uniform amplitude within their respective bandwidths and zero amplitude outside.

\begin{figure}[h!]
    \centering
    \includegraphics[width=0.7\textwidth]{b.png}
    \caption{Passband Signal}
    \label{fig:passband}
\end{figure}

Figure \ref{fig:passband} displays the frequency spectrum of the passband signal, which has been up-converted to a carrier frequency centered around 2.45 GHz. This spectrum spans a total bandwidth of 20 MHz, ranging from 2.44 GHz to 2.46 GHz, mirroring the bandwidth of the baseband signal but shifted to the carrier frequency. The rectangular shape signifies an ideal passband filter response, encapsulating the frequencies present in the transmitted signal.

The transformation from baseband to passband involves shifting the baseband frequencies to the carrier frequency, a process critical for RF transmission. This shift allows the signal to be transmitted over a medium such as air or cable at higher frequencies. The I and Q components in the baseband enable complex modulation schemes, facilitating efficient data transmission by utilizing both amplitude and phase variations.

\newpage
\section*{Problem \#3}

Consider a communication system with a very large $K$ factor for the Rician channel model. 

\begin{enumerate}
    \item Is it likely that there is a line-of-sight path or not? 
    \item Assume that the phase of the received signal increases by $\pi$ rad/s and the transmitter’s location is fixed. If the carrier frequency is 3 GHz, what is the speed of the receiver?
\end{enumerate}

\section*{Solution}

\subsection*{Part a}

In the Rician channel model, the $K$ factor represents the ratio of the power in the direct line-of-sight path to the power of the multipath components. A very large $K$ factor indicates that the power in the direct line-of-sight path is much greater than the power of the scattered paths. Therefore, if a communication system has a very large $K$ factor, it is likely that there is a strong line-of-sight path between the transmitter and the receiver.

\subsection*{Part b}

The phase change in the received signal due to the Doppler effect, when there is relative motion between the transmitter and receiver, is given by:

\[
f_d = \frac{\Delta \phi}{2\pi} \times f_s
\]

where:
\begin{itemize}
    \item $\Delta \phi$ is the phase change per second (rad/s),
    \item $f_s$ is the signal frequency (Hz).
\end{itemize}

The speed of the receiver $v$ can be calculated from the Doppler shift using the formula:

\[
f_d = \frac{v}{\lambda}
\]

where:
\begin{itemize}
    \item $\lambda$ is the wavelength of the carrier signal (m),
    \item $f_d$ is the Doppler shift (Hz).
\end{itemize}

The wavelength $\lambda$ can be found from the carrier frequency $f_c$ using:

\[
\lambda = \frac{c}{f_c}
\]

where:
\begin{itemize}
    \item $c$ is the speed of light ($3 \times 10^8$ m/s),
    \item $f_c$ is the carrier frequency (Hz).
\end{itemize}

Given:
\begin{itemize}
    \item $\Delta \phi = \pi$ rad/s,
    \item $f_c = 3$ GHz.
\end{itemize}

We can calculate the speed of the receiver $v$:

\[
f_d = \frac{\pi}{2\pi} = 0.5 \text{ Hz}, \quad \lambda = \frac{3 \times 10^8}{3 \times 10^9} = 0.1 \text{ m}, \quad v = f_d \cdot \lambda = 0.5 \times 0.1 = 0.05 \text{ m/s}.
\]

Therefore, the speed of the receiver is 0.05 m/s.
\newpage
\section*{Problem \#4}

Consider two communication systems, A and B, operating over a 10 MHz bandwidth. System A uses a 16-QAM modulation scheme achieving reliable communication at an $E_b/N_0$ of 10 dB. System B uses a QPSK modulation scheme, achieving reliability at an $E_b/N_0$ of 5 dB. Both systems aim for a data rate of 10 Mbps.

\subsection*{a. Spectral Efficiency}

The spectral efficiency $\eta$ is defined as the data rate $R$ in bits per second per bandwidth $B$ in hertz, expressed as:
\[
\eta = \frac{R}{B}.
\]

For both systems operating at 10 Mbps over a bandwidth of 10 MHz, the spectral efficiencies are:
\[
\eta_A = \eta_B = \frac{10 \text{ Mbps}}{10 \text{ MHz}} = 1 \text{ bps/Hz}.
\]

\subsection*{b. Gap to Shannon's Bound}

Shannon's capacity formula gives the maximum possible data rate $C$ for a given bandwidth $B$ and signal-to-noise ratio (SNR):
\[
C = B \log_2(1 + SNR).
\]

For System A:
\[
C_A = 10 \times \log_2(1 + 10^{10/10}) \approx 34.59 \text{ Mbps}.
\]

For System B:
\[
C_B = 10 \times \log_2(1 + 10^{5/10}) \approx 20.57 \text{ Mbps}.
\]

The theoretical gaps to Shannon’s bound, based on actual data rates, are:
\[
\text{Gap}_A = 34.59 \text{ Mbps} - 10 \text{ Mbps} = 24.59 \text{ Mbps},
\]
\[
\text{Gap}_B = 20.57 \text{ Mbps} - 10 \text{ Mbps} = 10.57 \text{ Mbps}.
\]

\subsection*{c. System Design Comparison}

System B is designed better for practical applications, as it achieves reliable communication at a lower $E_b/N_0$ of 5 dB compared to 10 dB for System A, indicating higher efficiency in noisy environments despite similar spectral efficiencies.


\end{document}
